\documentclass[11pt]{article}
\usepackage[margin=1in]{geometry}
\usepackage{amsmath,amssymb,amsthm}
\usepackage{graphicx}
\usepackage[hidelinks]{hyperref}

\newtheorem{theorem}{Theorem}
\newtheorem{lemma}{Lemma}
\newtheorem{remark}{Remark}

\newcommand{\N}{\mathbb{N}}
\newcommand{\F}{\mathcal{F}}
\newcommand{\K}{\mathbb{K}}
\newcommand{\T}{\mathbb{T}}
\newcommand{\supp}{\operatorname{supp}}
\newcommand{\Isom}{\operatorname{Isom}}
\newcommand{\Sym}{\operatorname{Sym}}

\title{A Threshold-Family Subcase of the Isometry Group Problem}
\author{run \texttt{fa\_banach\_001}, agent \texttt{agent\_lane\_02}}
\date{2026-06-14}

\begin{document}
\maketitle

\section*{Status}

This is a candidate partial result for the final questions in Brech--Ferenczi--Tcaciuc,
\emph{Isometries of combinatorial Banach spaces}, arXiv:1907.06815 \cite{BFT}.
It does not solve the general classification problem.  It gives a complete
classification for a natural class of regular families defined by a
nondecreasing threshold function.

\section*{Source Question}

In Section 5 of \cite{BFT}, after proving that isometries of combinatorial
spaces are signed permutations, the authors ask:

\begin{quote}
Given a regular family $\F$, what are the permutations of the basis which
induce an isometry of $X_\F$?

For which combinatorial spaces can we explicitly describe the group of its
isometries?
\end{quote}

The source crop below is from page 11 of the arXiv PDF.

\begin{center}
\includegraphics[width=\textwidth]{figures/open_problem_crop.png}
\end{center}

\section*{Idea of the Proof}

For the threshold families below, each integer $n$ carries a purely
combinatorial label: the largest cardinality of a member of the family that
contains $n$.  Monotonicity of the threshold function makes this largest
cardinality exactly $g(n)$.  Therefore any permutation preserving the family
must preserve every level set of $g$.  Conversely, if a permutation preserves
the levels of $g$, then moving a finite set can only move its minimum to a
point with threshold at least as large as the old one, so membership in the
family is preserved.

\section*{The Partial Result}

Let $g:\N\to\N$ be nondecreasing with $g(n)\geq 1$ for all $n$.  Define
\[
  \F_g=\{\varnothing\}\cup
  \{F\in[\N]^{<\omega}: F\neq\varnothing \text{ and } |F|\leq g(\min F)\}.
\]
Let
\[
  L_m=\{n\in\N: g(n)=m\}.
\]
Empty level sets are ignored in products below.  Let $\T_\K$ denote the
scalar sign group: $\{-1,1\}$ over the reals and the unit circle over the
complex numbers.

\begin{theorem}
\label{thm:threshold}
The family $\F_g$ is regular.  A permutation $\pi$ of $\N$ induces an isometry
of $X_{\F_g}$, equivalently of $X_{\F_g}^*$, if and only if
\[
  g(\pi(n))=g(n)\quad\text{for every }n\in\N.
\]
Consequently every surjective linear isometry of $X_{\F_g}$ or of
$X_{\F_g}^*$ is a signed permutation preserving the level sets $L_m$, and
\[
  \Isom(X_{\F_g})\cong \T_\K^{\N}\rtimes \prod_{m\geq 1}\Sym(L_m),
\]
with the same permutation part for $\Isom(X_{\F_g}^*)$.
\end{theorem}

\begin{proof}
First $\F_g$ is hereditary: if $E\subseteq F\in\F_g$ and $E\neq\varnothing$,
then $|E|\leq |F|\leq g(\min F)$, while $\min E\geq \min F$ and $g$ is
nondecreasing, so $|E|\leq g(\min E)$.  The spreading property is similar:
if $F=\{n_1<\cdots<n_k\}\in\F_g$ and $m_i\geq n_i$ with
$m_1<\cdots<m_k$, then $k=|F|\leq g(n_1)\leq g(m_1)$.
The family contains all singletons because $g\geq 1$.

To prove compactness, view $\F_g$ as a subset of $2^\N$.  Suppose
$F_j\in\F_g$ converges pointwise to $A\subseteq\N$.  If $A$ were infinite and
$a=\min A$, then for all large $j$ the set $F_j$ would contain $a$ and the
first $g(a)+1$ elements of $A$, and would contain no integer below $a$.
Hence $\min F_j=a$ and $|F_j|\geq g(a)+1$, contradicting
$F_j\in\F_g$.  Thus $A$ is finite.  If $A$ is nonempty and $a=\min A$, then
for all large $j$, all points of $A$ belong to $F_j$ and no point below $a$
belongs to $F_j$; consequently $\min F_j=a$ and
$|A|\leq |F_j|\leq g(a)$.  Hence $A\in\F_g$.  Therefore $\F_g$ is compact.

Now define the family-invariant number
\[
  h_{\F_g}(n)=\sup\{|F|:F\in\F_g,\ n\in F\}.
\]
Because $g$ is nondecreasing, any $F\in\F_g$ containing $n$ satisfies
\[
  |F|\leq g(\min F)\leq g(n).
\]
On the other hand, the set $\{n,n+1,\ldots,n+g(n)-1\}$ belongs to $\F_g$ and
contains $n$, so $h_{\F_g}(n)=g(n)$.

Let $\pi$ be a permutation of $\N$ with
$\pi(\F_g)=\F_g$.  Then cardinalities of family members containing a point
are preserved, so
\[
  g(\pi(n))=h_{\F_g}(\pi(n))=h_{\F_g}(n)=g(n).
\]
Thus every family automorphism preserves the level sets of $g$.

Conversely, suppose $g(\pi(n))=g(n)$ for all $n$.  Take a nonempty
$F\in\F_g$ and let $r=\min F$.  Let $s\in F$ be such that
$\pi(s)=\min \pi(F)$.  Since $s\geq r$ and $g$ is nondecreasing,
\[
  g(\min\pi(F))=g(\pi(s))=g(s)\geq g(r)\geq |F|=|\pi(F)|.
\]
Thus $\pi(F)\in\F_g$.  Applying the same argument to $\pi^{-1}$ gives
$\pi(\F_g)=\F_g$.

Brech--Ferenczi--Tcaciuc prove that for regular families, all isometries of
combinatorial spaces and their duals are signed permutations, and that a
permutation induces an isometry exactly when it sends the family to itself
\cite[Theorem 10 and the following corollary]{BFT}.  Combining that theorem
with the preceding family-automorphism classification gives the asserted
description of $\Isom(X_{\F_g})$ and $\Isom(X_{\F_g}^*)$.
\end{proof}

\section*{Relation to the Source Example}

The final example of \cite{BFT} treats the case where the threshold value is
constant on consecutive intervals and increases by one from one interval to
the next.  The theorem above allows arbitrary nondecreasing threshold
functions: jumps may skip values, plateaus may be finite or infinite, and the
constant case recovers the full symmetric permutation group.

\section*{Limitations and Novelty Check}

This is a partial result only.  It does not classify the isometry groups of
all combinatorial spaces, nor does it address regular families whose
membership is not determined by the pair $(|F|,\min F)$.

Local duplicate checks searched the run indexes for arXiv:1907.06815,
combinatorial Banach spaces, isometry groups, permutations, and threshold
families.  The source paper itself was searched for threshold-family and
level-set language.  No existing run packet or attempt was found for this
subcase.  No external literature search was used, so novelty confidence is
local rather than global.

\section*{Verification Notes}

No computational step is used.  The proof reduces the isometry classification
to the source theorem of \cite{BFT} and proves the remaining family
automorphism statement directly.  The two proof points most worth checking are
the compactness argument for $\F_g$ and the converse direction showing that a
$g$-level-preserving permutation sends every $F\in\F_g$ back into $\F_g$.

\begin{thebibliography}{9}
\bibitem{BFT}
C. Brech, V. Ferenczi, and A. Tcaciuc,
\emph{Isometries of combinatorial Banach spaces},
arXiv:1907.06815.
\end{thebibliography}

\end{document}
